\section{收敛性结果与完整推导}
本节重构 Wang 等人\cite{wang2026practical}的 Algorithm~1 与 Theorem~1。
为便于逐项检查，保留任意正常数下的有限步界，不以数值搜索得到的多项式常数代替证明。
结论针对以下单个矩阵变量的理论算法；第\ref{sec:implementation}节列出的工程修改需另外核对。

\subsection{假设、标量映射与常数}
令 $f:\mathbb R^{m\times n}\to\mathbb R$，$g_t=\nabla f(W_t)$，假设：
\begin{enumerate}
\item $f_*:=\inf_W f(W)>-\infty$，记 $\Delta=f(W_0)-f_*$。
\item $\|\nabla f(X)-\nabla f(Y)\|_F\le L\|X-Y\|_F$。
\item 在迭代历史 $\mathcal F_t$ 给定后，新的批梯度满足
 $\mathbb E[G_t\mid\mathcal F_t]=g_t$，
 $\mathbb E[\|G_t-g_t\|_F^2\mid\mathcal F_t]\le\sigma^2/b$。
 批内样本独立，批大小为 $b$。这里写成条件形式，以明确自适应迭代中使用的新样本与过去的关系。
\end{enumerate}

固定 $a=3.4445$、$b_p=-4.7750$、$c=2.0315$，定义
\begin{equation}
 p(x)=a+b_px+cx^2,\quad q(s)=s p(s^2),\quad
 q_0(s)=s,\quad q_{j+1}(s)=q(q_j(s)).
\end{equation}
注意 $b_p$ 是多项式系数，与批大小 $b$ 不同。
当 $s>0$ 时令 $\chi(s)=q_5(s)/s$，并以连续延拓定义 $\chi(0)=a^5$。
取
\begin{equation}\label{eq:constants}
 C_2=\min_{0\le s\le1}\chi(s),\quad
 C_1=\max_{0\le s\le1}\chi(s),\quad
 q_{\max}=\max_{0\le s\le1}q_5(s),\quad
 C_3=\min\{C_1,q_{\max}\sqrt{\min(m,n)}\}.
\end{equation}
这些常数有限且 $C_2>0$，理由如下。由于 $c>0$ 且 $4ac-b_p^2>0$，
$p(x)\ge a-b_p^2/(4c)>0$ 对所有实数 $x$ 成立。
因而
\[
 \chi(s)=\prod_{j=0}^{4}p(q_j(s)^2)
\]
在 $[0,1]$ 上为严格正的连续函数，包括 $s=0$ 的延拓值。
紧致性给出式\eqref{eq:constants}中的极值；这里不需要假设经验 NS 逐步逼近 $1$。
另外 $C_2\le q_5(1)\le q_{\max}$ 且 $C_2\le C_1$，故 $C_2\le C_3$。

\subsection{奇异值映射与方向对齐}
理论迭代使用
\begin{align}
 B_t&=\beta B_{t-1}+(1-\beta)G_t, & B_{-1}&=0,\\
 M_t&=\beta B_t+(1-\beta)G_t, &
 \tau_t&=\max\{\|M_t\|_F,\varepsilon\},\\
 X_{t,0}&=M_t/\tau_t,& X_{t,j+1}&=p(X_{t,j}X_{t,j}^{\mathsf T})X_{t,j},\quad 0\le j<5,\\
 D_t&=X_{t,5},& W_{t+1}&=(1-\eta\lambda)W_t-\eta D_t.
\end{align}
若 $M_t=U\operatorname{diag}(\sigma_i)V^{\mathsf T}$ 是紧致 SVD，令 $s_i=\sigma_i/\tau_t$。
直接代入矩阵多项式并归纳可得
\[
 X_{t,j}=U\operatorname{diag}(q_j(s_i))V^{\mathsf T},\qquad 0<s_i\le1.
\]
零矩阵的输出为零；秩亏时只保留正奇异值，无须补齐零奇异值方向。
由 $C_2s\le q_5(s)\le C_1s$，
\begin{align}
 \langle M_t,D_t\rangle_F
 &=\tau_t\sum_i s_iq_5(s_i)
 \ge C_2\|M_t\|_F^2/\tau_t
 \ge C_2\|M_t\|_F-C_2\varepsilon/4,\label{eq:alignment}\\
 \|D_t\|_F&\le\min\left\{C_1\sqrt{\sum_i s_i^2},q_{\max}\sqrt{\operatorname{rank}(M_t)}\right\}
 \le C_3.\label{eq:direction}
\end{align}
最后一个对齐不等式分两种情况：当 $z=\|M_t\|_F\ge\varepsilon$ 时 $z^2/\tau_t=z$；
当 $0\le z<\varepsilon$ 时，$z^2/\varepsilon-z+\varepsilon/4=(z-\varepsilon/2)^2/\varepsilon\ge0$。
这样得到的界直接控制下降方向，不依赖 $\|D_t-UV^{\mathsf T}\|$ 随 NS 次数单调减小。

\subsection{主定理的有限步形式}
\begin{theorem}[文献 Theorem 1 的非优化常数形式]\label{thm:main}
在上述假设下，取常数 $c_1,c_2,c_3>0$ 和 $T\ge\max\{1,c_2^2\}$，设置
\begin{equation}\label{eq:parameters}
 \eta=c_1T^{-3/4},\quad 1-\beta=c_2T^{-1/2},\quad
 0<\varepsilon\le c_3T^{-1/4},\quad
 \lambda=\frac{C_2}{2(\|W_0\|_F+C_3\eta T)}.
\end{equation}
记 $C_4=C_2/2+C_3$、$A=2(C_2+C_3)/C_2$ 和 $E_0=\|B_{-1}-g_0\|_F$，则
\begin{align}\label{eq:finite}
 \frac1T\sum_{t=0}^{T-1}\mathbb E\|g_t\|_F
 \le{}&\left[\frac{2\Delta}{C_2c_1}
 +\frac{ALc_1C_4}{c_2}
 +\frac{A\sqrt{2c_2}\sigma}{\sqrt b}+\frac{c_3}{2}\right]T^{-1/4}\nonumber\\
 &+\frac{AE_0}{c_2}T^{-1/2}
 +\frac{Lc_1C_4^2}{C_2}T^{-3/4}.
\end{align}
因此，固定问题与常数后，平均期望梯度 Frobenius 范数为 $O(T^{-1/4})$。
\end{theorem}

\subsection{权重、总更新和动量误差的控制}
先证明权重衰减不会破坏所需的方向界。由 $C_2\le C_3$ 以及式\eqref{eq:parameters}，
$0<\eta\lambda\le1/(2T)<1$。由式\eqref{eq:direction}归纳得到
\begin{equation}\label{eq:weight}
 \|W_t\|_F\le\|W_0\|_F+\eta C_3t,\quad
 \lambda\|W_t\|_F\le C_2/2,\quad
 \|H_t\|_F\le C_4,\qquad H_t=D_t+\lambda W_t.
\end{equation}
其中前一个不等式来自
$\|W_{t+1}\|_F\le(1-\eta\lambda)\|W_t\|_F+\eta C_3\le\|W_t\|_F+\eta C_3$。

令 $e_t=M_t-g_t$，$\zeta_t=G_t-g_t$。展开两个动量递推，得到
\begin{align}
 M_t&=\beta^{t+2}B_{-1}
 +(1-\beta)\sum_{j=0}^{t-1}\beta^{t-j+1}G_j+(1-\beta^2)G_t,\\
 e_t&=\beta^{t+2}(B_{-1}-g_0)
 -\sum_{j=1}^{t}\beta^{t-j+2}(g_j-g_{j-1})+N_t,\label{eq:tracking-decomposition}\\
 N_t&=(1-\beta)\sum_{j=0}^{t-1}\beta^{t-j+1}\zeta_j+(1-\beta^2)\zeta_t.
\end{align}
式\eqref{eq:tracking-decomposition}可通过把相邻梯度差的和展开、抵消内部项核查。
光滑性与式\eqref{eq:weight}给出
$\|g_j-g_{j-1}\|_F\le L\eta C_4$。
条件无偏性使不同时间噪声的交叉内积期望为零，故
\begin{align}
 \mathbb E\|N_t\|_F^2
 &\le\frac{\sigma^2}{b}\left[(1-\beta^2)^2
 +(1-\beta)^2\sum_{j=0}^{t-1}\beta^{2(t-j+1)}\right]\\
 &\le\frac{\sigma^2}{b}(1-\beta)\frac{1+2\beta-2\beta^3}{1+\beta}
 \le\frac{2(1-\beta)\sigma^2}{b}.
\end{align}
再用三角不等式、几何级数和 Cauchy--Schwarz 不等式，得到
\begin{equation}\label{eq:tracking-bound}
 \overline E_T:=\frac1T\sum_{t=0}^{T-1}\mathbb E\|e_t\|_F
 \le\frac{E_0}{T(1-\beta)}
 +\frac{L\eta C_4}{1-\beta}
 +\sigma\sqrt{\frac{2(1-\beta)}{b}}.
\end{equation}
初值项与梯度漂移项各使用一次几何和上界 $1/(1-\beta)$；噪声项使用二阶矩界后开平方。
以上推导也覆盖 $\beta=0$，此时空和按零处理。

\subsection{下降不等式、求和与收敛指标}
\begin{proof}[定理\ref{thm:main}的证明]
光滑性给出
\[
 f(W_{t+1})\le f(W_t)-\eta\langle g_t,H_t\rangle_F
 +\frac{L\eta^2}{2}\|H_t\|_F^2.
\]
把线性项拆开，利用式\eqref{eq:alignment}和\eqref{eq:weight}：
\begin{align}
 \langle g_t,H_t\rangle_F
 &=\langle M_t,D_t\rangle_F-\langle e_t,D_t\rangle_F
   +\lambda\langle g_t,W_t\rangle_F\\
 &\ge C_2\|M_t\|_F-C_2\varepsilon/4-C_3\|e_t\|_F
   -\lambda\|W_t\|_F\|g_t\|_F\\
 &\ge\frac{C_2}{2}\|g_t\|_F-(C_2+C_3)\|e_t\|_F-C_2\varepsilon/4.
\end{align}
因此，取期望后对 $t=0,\ldots,T-1$ 求和，以 $f(W_T)\ge f_*$ 消去终点，得到
\begin{equation}\label{eq:pre-rate}
 \frac1T\sum_{t=0}^{T-1}\mathbb E\|g_t\|_F
 \le\frac{2\Delta}{C_2\eta T}+A\overline E_T
 +\frac\varepsilon2+\frac{L\eta C_4^2}{C_2}.
\end{equation}
把式\eqref{eq:tracking-bound}及参数选择代入，三项分别为
\[
 \frac{E_0}{c_2}T^{-1/2},\qquad
 \frac{Lc_1C_4}{c_2}T^{-1/4},\qquad
 \frac{\sqrt{2c_2}\sigma}{\sqrt b}T^{-1/4}.
\]
其余项直接得到式\eqref{eq:finite}，证明完成。
\end{proof}

若 $\sigma>0$ 且 $\Delta>0$，可进一步平衡领先系数。忽略可自由选小的 $c_3/2$ 后，
先对 $c_1$ 极小化，得到
$c_1^2=\Delta c_2/[(C_2+C_3)LC_4]$；再平衡 $c_2^{-1/2}$ 与 $c_2^{1/2}$ 项，得到
\[
 c_2=\sqrt{\frac{2C_4L\Delta b}{(C_2+C_3)\sigma^2}}.
\]
于是领先尺度为 $(L\Delta\sigma^2/b)^{1/4}$ 乘上 NS 常数，与文献主结果的量级一致。
若 $\sigma=0$ 或 $\Delta=0$，不使用含除零的优化常数公式，保留式\eqref{eq:finite}及任意正常数即可。

定理控制的是驻点指标：均匀随机选取 $R\in\{0,\ldots,T-1\}$ 时，
$\mathbb E\|\nabla f(W_R)\|_F$ 等于上述平均值。
它不保证每一步梯度单调下降，也不保证非凸问题达到全局最优。
令 $\lambda=0$，同样推导中取 $C_4=C_3$ 并去除衰减项，可得到无权重衰减的对应结论。
